\scanpage{532}
\mspage{443}
% archive sheet number: 531

\uline{Corollaire} 4.15.7.3. Supposons que $I$ soit grd devant $\pi_0$, où $\pi_0$ est un cardinal infini tel que $\pi_0\geq\operatorname{card}\deleted{\unsure{\mathrm{red}|\mathrm{nerf}|\mathrm{morph.}}} C$. Alors $(x_i)_{i\in I}$ est ess.t constant sss on peut trouver un foncteur cofinal $J\longrightarrow I$, avec $J$ \deleted{\unsure{ordonné et}} grand devant $\pi_0$\inserted{(et \unsure{réduit si on veut|réordonné})}, de telle façon que le système réindexé $(x_j)_{j\in J}$ satisfasse les conditions suivantes, \deleted{il existe $1)$ $\forall j\in\mathrm{Ob}\,J$}

% slanted
\margin{La démonstration \unsure{légèrement réarrangée|convenablement réarrangée} montrera qu'on peut prendre pour $J$ une sous-catégorie (pas pleine) d'une des catégories $i_{\unsure{1|0}}\backslash I$, pour un objet convenable $i_1$ \deleted{\unsure{$j$}} de $I$. \deleted{\illegible{1}} On a donc $\operatorname{card}J\leq\pi$ si $\operatorname{card}I\leq\pi$, $\pi$ étant un cardinal infini donné. \deleted{\unsure{Lorsque $I$ admet} \illegible{2}}}

\deleted{$x_j=\unsure{x_0}$} $1^\circ)$ L'application $j\mapsto x_j$ de $\mathrm{Ob}\,J$ dans $\mathrm{Ob}\,C$ est constante, i.e. il existe un $x_0\in\mathrm{Ob}\,C$ tel que $x_j=x_0\quad\forall j\in\mathrm{Ob}\,J$ ; $2^\circ)$ il existe un projecteur $p$ dans $x_0$ tel que pour $j<j'$ dans $J$, on ait $\tau_{j',j}=p$ ; $3^\circ)$ Le projecteur $p$ admet une image dans $C$.

\uline{Dém} : Ces conditions impliquent \inserted{évidemment} celles de 4.15.7.1, en prenant $p_j=p\quad\forall j\in\mathrm{Ob}\,J$, du moins si $J$ n'admet pas de plus grd élément (ce qui fait une difficulté pour $2^\circ)$). Mais si $J$ admet un objet final, \inserted{\unsure{il en est de même de $I$, et|il en est de même de $I$, it}} il est clair \inserted{alors} que \deleted{\unsure{tt}} syst. inductif indexé par $\deleted{J}I$ \deleted{\unsure{admet}} est ess. constant. Donc, dans \deleted{les deux} tous les cas, les conditions $1^\circ)$ $2^\circ)$ $3^\circ)$ ci-dessus impliquent que $\underline{x}$ est ess. constant. Inverse%
